% Abstract (Task 46): the draft of handoff/42_srep_plan.md section 5, unstructured, no citations, at most 200 words.
\begin{abstract}
Attention weights of transformer-based network intrusion detectors are often read as explanations, and agreement
between attention and attribution methods has been proposed as a signal of attack.
We tested both readings on feature-tokenizer transformers trained on LYCOS-IDS2017 and LycoS-Unicas-IDS2018, with
\nSeeds{} seeds and near-duplicate rows removed from evaluation.
Rollout attention agreed weakly to moderately with attribution methods, was less faithful than integrated gradients
and occlusion, and changed little under protocol-constrained attacks that targeted it.
Consistency scores added to a confidence-and-distance detector helped against an attacker that ignored them, as did
controls without attention.
Against attackers that optimised the exact score of each detector, with surrogate gradients only for \nRankScores{}
rank-based scores, the consistency scores raised the share of successful and undetected attack rows by
\axHeadVoneSeventeenAbs{} on 2017 (\paramCILevel{} confidence interval \axHeadVoneSeventeenCI) and by
\axHeadVoneEighteenAbs{} on 2018 (\axHeadVoneEighteenCI), with \axHeadOppSeedsPhrase{} of the opposite sign.
An attacker that used differentiable stand-ins for these scores had shown a reduction of \axStandinVoneSeventeenAbs{}
on 2017.
Giving the attacker of the base detector the same iterate selection changed the outcome of \sxOutcomeChangedPhrase.
Adaptive evaluations of explanation-based attack detectors need attackers that optimise the exact detection score.
\end{abstract}
